%!TEX TS-Program = xelatex

\documentclass{beamer}

%\usepackage[latin2]{inputenc}
%\usepackage[T1]{fontenc}
%\usepackage[polish]{babel}
\usepackage{polski}
\usepackage{graphicx}
\usepackage{minted}
%\usepackage{epsfig}
\usepackage{xcolor}
\usetheme{metropolis}
%\usepackage{minted}
%\usetheme{JM}
\usepackage{hyperref}
\usepackage{textcomp}
\usepackage{xcolor}
\usecolortheme{seahorse}


\usepackage{color}
\usepackage{listings}

\newtheorem{definicja}{Definicja}
\newtheorem{twierdzenie}{Twierdzenie}
\newtheorem{przyklad}{Przykład}
\newtheorem{uwaga}{Uwaga}
\newtheorem{algorytm}{Algorytm}

\author{}
\title{Gnuplot \\ Tworzenie wykresów}
\date{2025}

\begin{document}
\maketitle

\section{Wstęp}

% \frame{{Narzędzia alternatywne}

% \begin{columns}
% \begin{column}{0.5\textwidth}
% \begin{itemize}
% \item Fityk
% \item Grace
% \item Veusz
% \item SciDAVis
% \item QtiPlot
% \item PAW++
% \item Root (podstawowy w fizyce wysokich energii)
% \item Matlab
% \item Octave
% 
% \end{itemize}
% \end{column}
% \begin{column}{0.5\textwidth}
% \begin{itemize}
% \item Maple
% \item Mathematica
% \item XGraph
% \item SciPy, NumPy
% \item Scilab
% \item R
% \item Origin
% \item Excel
% \item ....
% \end{itemize}
% \end{column}
% \end{columns}
% }
% \frame{
% \begin{itemize}
%  \item Gnuplot jest jednym z najbardziej popularnych programów do tworzenia ,,grafiki naukowej'', czyli wykresów. 
%  \item Występuje on w wersjach pracujących na większości systemów.
%  \item Zaletami są efekty jego pracy i pełna kontrola nad wynikami. 
%  \item Niestety jest to program wymagający znajomości odpowiednich instrukcji.
%  \item Obecna wersja stabilna to 6.0.1 (listopad 2024) \\
%  \pause
%  
% Najnowsze wersje skupiają się na usprawnieniach i usuwaniu znalezionych błędów.\\
% Taka mrówcza ale niezwykle ważna praca.
% \end{itemize}
% 
% 
% }
%\frame{{Reklama -- przegląd możliwości}
% \begin{itemize}
% \item Rysowanie wykresów funkcji i zbiorów danych.
% \item Rysowanie obiektów 2- i 3-wymiarowych.
% \item Współrzędne kartezjańskie, sferyczne, biegunowe, parametryczne.
% \item Dopasowanie funkcji do danych -- także funkcje nieliniowe wraz z dokładnością oszacowania parametrów (!).
% \item Transformacje danych.
% \item Generowanie wykresów w różnych formatach graficznych \\
% w szczególności w postscripcie (grafika wektorowa).
% \item Generowanie wykresów złożonych \\
% wstawki, różne skale, wiele wykresów na jednym ...
% \item \alert{Możliwe jest przygotowanie skryptu generującego wykres -- automatyzacja pracy.}
% \item Animacje
% \end{itemize}
% }
% %\frame{{Trochę przykładów}

% \includegraphics<1>[width=8cm]{przykl_1.eps} 
% %\onslide<2> 
% %Przykład złożony \\
% \includegraphics<2>[width=8cm]{przykl_2.eps} 
% %\onslide<3>
% %Wykres konturowy \\
% \includegraphics<3>[width=8cm]{przykl_3.eps} 

% %\onslide<4>
% %Wykres konturowy \\
% \includegraphics<4>[width=8cm]{przykl_4.eps} 

% %\onslide<5>
% %Wykres parametryczny \\
% \includegraphics<5>[width=8cm]{przykl_5.eps} 

% %\onslide<6>
% %Układ biegunowy \\
% \includegraphics<6>[width=8cm]{przykl_6.eps} 

% %\onslide<7>
% %Aproksymacje \\
% \includegraphics<7>[width=8cm]{przykl_7.eps} 

% %\onslide<8>
% %3D \\
% \includegraphics<8>[width=7cm]{przykl_8.eps} 

% %\onslide<9>
% %Kolory \\
% \includegraphics<9>[width=8cm]{przykl_9.eps} 

% %\onslide<10>
% %pm3D \\
% \includegraphics<10>[width=8cm]{przykl_10.eps} 

% }

%\frame{{Jak działa? Gdzie zdobyć? Literatura.}

% \begin{itemize}
% \item<1-> Sterowanie odbywa się przez \alert{konsolę} lub odpowiednie pliki wsadowe.
% \item<2-> Program i aktualne instrukcje można pobrać ze strony: \\
% http://www.gnuplot.info/
% \item<3-> Użyteczne odnośniki: \\
% http://www.fuw.edu.pl/\textasciitilde pablo/pracownia/gnuplot.pdf \\
% http://www.knf.ifd.uni.wroc.pl/materialy/gnuplot.pdf
% \item<4-> W programie -- przez słowo kluczowe {\color{green} help \it komenda} lub {\color{green} ? \it komenda}
% \item<5-> W dokumentacji programu:\\
%  gpcard.pdf\\
%  gnuplot.$\ast$; $\ast\Rightarrow $ html, pdf, info, hlp ...\\
% pliki demonstracyjne -- katalog demo; 
% \item<6-> Podręczniki: 
% \begin{itemize}
%  \item Philipp K. Janert: Gnuplot in Action (2016)
%  \item Lee Philips: Gnuplot 5 (2017)
% \end{itemize}

% \end{itemize}
% }
\frame{{Wprowadzenie }
Praca z gnuplotem może być wykonywana
\begin{enumerate}
\item w trybie interaktywnym {\color{red} (nie polecam) }
\item oparta na tworzeniu plików wsadowych {\color{red} (zalecana)}.
\end{enumerate}
Ta druga metoda ma tą  zaletę, że pozwala na łatwe modyfikowanie plików wsadowych i dopasowywanie ich do swoich potrzeb. \\
\begin{itemize}
 \item Zwykle mamy do czynienia z powtarzalnymi czynnościami. 
 \item Dodatkową zaletą jest możliwość przygotowania (użycia) skryptów automatyzujących pracę.
 \item Łatwość zachowania stylu (formatu) przygotowywanych wykresów.
\end{itemize}


}
\frame{{Podstawowe instrukcje }
Program uruchamiamy wpisując z linii poleceń (w terminalu) instrukcje:
\begin{enumerate}
\item dla sesji interaktywnej -- {\color{blue} gnuplot} - instrukcje są wpisywane zgodnie z potrzebami
\item dla sesji opartej na pliku wsadowym -- {\color{blue} gnuplot plik1 plik2} - instrukcje znajdują się w plikach
\item względnie przygotowujemy skrypt wykonywalny, który zawiera wywołanie programu jak i instrukcje do wykonania.
\end{enumerate}
}
\frame{{Własności konsoli}
 \begin{itemize}
  \item Rozróżnia litery małe i duże.
\item Polecenia posiadają (zwykle) opcjonalne parametry.
\item Opis poleceń jest dostępny pod komendą: ``help''.
\item Istnieją skróty poleceń np. {\color{blue} plot $\Rightarrow$ p}.
\item Polecenia mogą zająć kilka wierszy, kontynuację w kolejnym wierszu umożliwia {\color{blue} $\backslash$} na końcu wiersza.
\item Napisy umieszcza się pomiędzy znakami apostrofu lub cudzysłowu.
\item Historię poleceń przewijamy $\uparrow \downarrow$.
 \end{itemize}
}
% \frame{{Instrukcje edytorskie: }
%     
% \begin{minipage}[t]{5cm}
% \begin{description}
% \item[\color{blue} \^{} B \color{black}] --  wstecz o jedną pozycję
% \item[\color{blue} \^{} F \color{black}] --  naprzód o jedną pozycję 
% \item[\color{blue} \^{} A \color{black}]  -- na początek linii 
% \item[\color{blue} \^{} E \color{black}] -- na koniec linii 
% \item[\color{blue} \^{} H , DEL\color{black}] -- usuwa poprzedni znak 
% \item[\color{blue} \^{} D \color{black}] -- usuwa bieżący znak 
% \end{description}
% \end{minipage}
% \begin{minipage}[t]{5.5cm}
% \begin{description}
% \item[\color{blue} \^{} K \color{black}] --  usuwa od pozycji do początku linii
% \item[\color{blue} \^{} U \color{black}] --  usuwa całą linię
% \item[\color{blue} \^{} W \color{black}]  -- usuwa od pozycji do końca linii
% \item[\color{blue} \^{} P \color{black}] -- przesuwa po historii poleceń wstecz
% \item[\color{blue} \^{} N \color{black}] -- przesuwa po historii poleceń do przodu
% \end{description}
% \end{minipage}
% }
\section{Wykresy}
\frame{{Wykresy}
Dostępne są 3 instrukcje generujące wykresy:
\begin{description}
\item[{\color{blue} plot (p) }] generuje 2-wym wykres
\item[{\color{blue} splot (sp) }] generuje 3-wym wykres
\item[{\color{blue} replot }] powtarza rysowanie wykresów, np. po wprowadzeniu dodatkowych zmiennych
\end{description}
Te same polecenia służą do sporządzania wykresu danych i funkcji.
Instrukcje {\color{blue} plot i splot} wymagają podania danych lub funkcji do wykresu oraz sposobu wykonania wykresu -- stylu punktów lub linii.
}

\frame{{Wykresy -- definiowanie danych }
Plik danych dla gnuplota powinien być utworzony w postaci kolumny liczb (miejsca dziesiętne oddzielamy {\color{red} kropką}!). 

\vspace{1cm}
Komentarze co do zawartości pliku muszą być poprzedzone znakiem:\\ \begin{center}
\#
\end{center}  
linia rozpoczynająca się od tego znaku jest ignorowana. W ten sam sposób możemy pisać komentarze w plikach wsadowych. 
}

\frame{{Wykresy -- definiowanie danych }
Domyślnie pierwsza kolumna jest traktowana jako wartości dla "`x"' \\ 
a druga dla "`y"'. 

\vspace{1cm}
W przypadku bardziej złożonych plików możliwy jest wybór odpowiednich kolumn poprzez instrukcje {\color{blue} using n:k} lub {\color{blue} u n:k} gdzie "`n"' i "`k"' są numerami kolumn nas interesujących. 

\vspace{1cm}
Podobnie dla przypadku 3 wym.: {\color{blue} using n:k:l} lub {\color{blue} u n:k:l}}

\frame{{Wyrażenia arytmetyczne}
\begin{itemize}
 \item Liczby zespolone: \{1,0\}.
\item Operator potęgowania: $\ast \ast$.
\item Operator ``silnia'': !.
\item ``Numer kolumny'': \$, np: p 'x.txt' u (\$1 + \$4):(\$3 **4). Co odpowiada w układzie parametrycznym $(x=x_1+x_4, y= x_3^4$).
\item Pozostałe operatory jak w języku C: !, {\~{} }, \&, |, $\ast$, \% (reszta z dzielenia), /, +, -, ==, !=, <, <=, >, >=, \&\& (i), || (lub).
\item Operator ? :, np. plot x < 0 ? -x : x.
\end{itemize}
}

\frame{{Funkcje }
Gnuplot posiada bogaty zbiór zdefiniowanych funkcji. \\
{\color{blue} $abs(x)$} - wartość bezwzględna (też dla liczb zespolonych),\\
{\color{blue} $acos(x)$} -  $\arccos$ - wartość w radianach, bądź stopniach, w zależności od ustawienia parametru globalnego {\color{blue} angles} \\
{\color{blue} $acosh(x)$} - $arccosh$, {\color{blue} $asin(x)$} - $\arcsin(x)$,\\
 {\color{blue} $asinh(x)$} -  $arsinh(x)$ - w radianach \\ 
{\color{blue} $arg(x)$} - zwraca fazę liczby zespolonej (w radianach lub stopniach) \\
%}
%\frame{{Funkcje}
{\color{blue} $ atan(x) $} - $\arctan(x)$ - w radianach lub stopniach  \\ 
{\color{blue} $ atan2(y,x)$} - $\arctan(Re(\frac{y}{x}))$ - w radianach lub stopniach \\ 
{\color{blue} $ atanh(x)$} - $ \arctan(x)$ - w radianach\\ 
{\color{blue} $ceil(x) $}, {\color{blue} $floor(x) $}  - część całkowita \\ 
{\color{blue} $\cos(x) $, $\cosh(x) $}, {\color{blue} $ \sin(x) , \sinh(x)$} , {\color{blue} $\tan(x), \tanh(x) $}, {\color{blue} $\exp(x) $}   \\ 
}
\frame{{Funkcje}

{\color{blue} $\log(x) $} - $\log(x)$, $log10(x)$ - $\log_{10}(x)$ \\ 
{\color{blue} $norm(x) $} -  wartość rozkładu normalnego \\ 
{\color{blue} $rand(x) $} - liczba pseudolosowa z przedziału $ \left[ 0,1 \right] $  \\ 
{\color{blue} $ sgn(x)$} -  $+1 $ gdy $x>0$, $-1$ gdy $x<0$, \\
{\color{blue} $sqrt(x) $} - $\sqrt{x}$  \\ 
 i kilka innych ...
}


\frame{{Funkcje i zmienne}
Funkcje definiujemy wykorzystując wbudowane operacje:\\
{\color{blue}
f(x)=a*sin(x)+x**3, czyli $f(x)=a\sin(x) +x^3$\\ 
g(x,y)=x**2 +y**3 \\
k(x) = (x != 0) ? sin(x)/x : 1 }\\
Natomiast zmiennym przypisujemy wartość: \\
a=2 \\
Tak zdefiniowaną funkcję możemy narysować \\
{\color{blue} p f(x) w l} \\
Istniejącą listę funkcji użytkownika możemy uzyskać poprzez polecenie: {\color{blue} show functions}\\
oraz listę zmiennych {\color{blue} show variable}\\
liczba $\pi$ jest zmienną predefiniowaną: {\color{blue} pi}.
}

\frame{{ Wykresy -- linie}
Dostępnych jest bardzo wiele stylów linii: 
\begin{figure}
	\centering
	\includegraphics[angle=-90,scale=0.35]{test.eps}
% test.eps: 300dpi, width=4.27cm, height=6.10cm, bb=50 50 554 770
%	\caption{Style grafiki i linii}
	\label{fig:linie}
\end{figure}

}


\frame{{Wykresy -- linie }

Zatem instrukcja, która wygeneruje wykres na podstawie pliku danych ma postać \\
{\color{blue} p "`dane.dat"' w l } \\
aby dodatkowo zdefiniować rodzaj lini należy zastosować: \\
{\color{blue} lt (nr)} - rodzaj linii i  {\color{blue} lw (nr)} - grubość linii \\ zatem instrukcja nasza ma postać:
{\color{blue} p "`dane.dat"' w l lt 1 lw 2}  \\
Zamiast wpisywać te wartości za każdym razem można zdefiniować interesujący nas styl linii:\\
{\color{blue} set linestyle 1 lt 2 lw 2 } \\
Wtedy odwołanie będzie miało postać:
{\color{blue} plot sin(x) ls 1} \\

}

% \frame{{Wykresy -- linie}
% \begin{minipage}[t]{5cm}
% \begin{figure}
% 	\centering
% 	\includegraphics[width=5cm]{przyklad_1.eps}
% 	\caption{przykład 1}
% \end{figure}
% \end{minipage}
% \begin{minipage}[t]{5cm}
% \begin{figure}
% 	\centering
% 	\includegraphics[width=5cm]{przyklad_2.eps}
% 	\caption{przykład 2}
% \end{figure}
% \end{minipage}
% }


\frame{{Wykresy -- punkty }
Podobnie można zrobić dla punktów. \\
{\color{blue} p "`dane.dat"' w p } \\
Definicje rodzaju i wielkości punktów otrzymujemy wpisując: \\
{\color{blue} pt (nr)} - rodzaj punktu i  {\color{blue} ps (nr)} - wielkość punktów \\ zatem instrukcja nasza ma postać: \\
{\color{blue} p "`dane.dat"' w p pt 1 ps 1}  \\
Uwaga należy stosować umiarkowane wartości parametru {\color{blue} ps}


}

\frame{{Wykresy -- punkty}
\begin{minipage}[t]{5cm}
\begin{figure}
	\centering
	\includegraphics[width=5cm]{przyklad_3.eps}
	\caption{przykład 3}
\end{figure}
\end{minipage}
\begin{minipage}[t]{5cm}
\begin{figure}
	\centering
	\includegraphics[width=5cm]{przyklad_4.eps}
	\caption{przykład 4}
\end{figure}
\end{minipage}
}
\begin{frame}
\frametitle{Style linii}

Dostępne rodzaje prezentacji danych:

\begin{tabular}{lll}
\hline
lines & points & linespoints \\
impulses & dots & steps \\
fsteps & histeps & errorbars \\
xerrorbars & yerrorbars & xyerrorbars \\
boxes & boxerrorbars & boxxyerrorbars \\
financebars & candlesticks & vector \\
\hline
\end{tabular}

\end{frame}



\frame{{Zakresy}
Zakres zmiennej podajemy wykorzystując nawiasy kwadratowe {\color{blue} []}. \\
Nawiasy te mogą się odnosić do każdej z osi wykresu z zachowaniem kolejności:\\
$[x_1:x_2] [y_1:y_2] [z_1:z_2]$ \\
Zakresy zmiennych podajemy wykorzystując ``:'', z tym że możliwe jest:\\
\begin{itemize}
 \item podanie pełnego zakresu zmiennych: [-3:5]
\item domyślny początek lub koniec zmiennych: [2:] lub [:4]
\item określenie zakresu tylko jednej z osi: [:] [3:] [2:4]
\end{itemize}
}

\frame{{Wiele wykresów}
Wykres kilku funkcji:\\
{\color{blue} p [0:2*pi] x w l, sin(x) w l, cos(x) w l  } \\

W jednej linii oddzielone przecinkami wypisujemy funkcje, które mają być na tym samym wykresie.

}

\frame{{Próbkowanie}
\begin{itemize}
 \item Wykresy są rysowane linią łamaną z {\color{blue}   ``samples'' } wierzchołkami. \\
\alert{ Sprawdzić jak wygląda wykres funkcji tangens!} 
\item Wartość tego parametru możemy sprawdzić: {\color{blue}   ``show samples'' }
\item Zmienić ustawienia: {\color{blue}   ``set samples 500'' }
\item Zmienna ta jest także wykorzystywana przy interpolacji.
\end{itemize}
}

\frame{{Wykresy biegunowe i parametryczne}
\begin{itemize}
 \item Współrzędne biegunowe: \\
 {\color{blue}   set polar \\
 p [0:4*pi] t  \\
 p [0:4*pi] t, 2*t }
 \item Wykresy parametryczne \\
 {\color{blue}   set parametric \\
 p [0:2*pi] sin(t), cos(t)  \\
 p [0:4*pi] t*sin(t), t*cos(t) }
 \item Powrót do współrzędnych kartezjańskich zapewnia: \\
 {\color{blue}   unset polar \\
unset parametric  }
\item Domyślną zmienną zależną jest x dla wykresów kartezjańskich \\ y = f(x) i kąt t dla wykresów biegunowych, r = f(t).  \\ Dla wykresów parametrycznych niezależne jest t 
natomiast rysowane funkcje to x(t) i y(t) -- para funkcji.
\end{itemize}

}

\frame{{Skala logarytmiczna}
Zależności potęgowe lub wykładnicze najwygodniej jest obserwować w skalach logarytmicznej lub półlogarytmicznej $\rightarrow$ otrzymujemy wykres prostej.\\
\begin{itemize}
 \item Logarytmiczna oś ``x'', podstawa 10:\\
{\color{blue} set logscale x}
\item Logarytmiczna oś ``y'', podstawa 2:
{\color{blue} set logscale y 2 }
\item Logarytmiczne osie ``x'' i ``y'':
{\color{blue} set logscale xy } 
\item Powrót do współrzędnych kartezjańskich:
{\color{blue} unset logscale xy}
\end{itemize}
}

\frame{{Wykresy danych}
Opcje dla funkcji w większości są dostępne przy wykonywaniu wykresów zbiorów danych.
\begin{itemize}
 \item Pliki z danymi mogą mieć dowolną liczbę kolumn: \\
{\color{blue} p 'dane.dat' u 7:2 }
\item Pliki z danymi mogą zawierać dowolną liczbę "`bloków"' danych;
\item Bloki oddziela się od siebie pustymi wierszami. 
\item Do bloków odwołujemy się opcją index: \\
{\color{blue} p 'dane.dat' index 2:4} 
\item Numeracja bloków zaczyna się od 0. 
\item Wykres bloków 0,3,6,9,12: \\
{\color{blue} p 'dane.dat' index 0:12:3}
\item Do samodzielnego sprawdzenia opcja {\color{blue} every}.
\end{itemize}
}

\frame{{Manipulowanie danymi -- {\color{blue} using}}
\begin{itemize}
 \item  Wykres kolumny 3 względem 1: {\color{blue} p [:] 'dane.dat' u 1:3 }
\item Wykres kolumny 2 względem x = 0, 1, ...: {\color{blue} p 'dane.dat' u 0:1}
\item  Wykres sumy kolumn 2 i 3 względem kol. 1: \\ 
{\color{blue} p 'dane.dat' u 1:(\$2 + \$3)} \\
\$n oznacza liczbę odczytaną z n-tej kolumny
\item Transformacja danych: {\color{blue} p 'dane.dat' u 1:(log(abs(\$2)))}
\item Filtrowanie danych: {\color{blue} p 'dane.dat' u 1:(\$2 < 10 ? \$2 : 1/0)}
\end{itemize}
}

\frame{{Interpolacja i ekstrapolacja}
\begin{itemize}
 \item Aproksymacja danych krzywą ciągłą: \\
{\color{blue} plot 'dane.dat' u 1:2:(1.0) smooth acsplines }
\item Trzeci parametr opcji {\color{blue} using} określa stopień wygładzania danych; im jest mniejszy, tym wykres będzie słabiej "`przyciągany"' do punktów wykresu.
\item Interpolacja danych krzywą ciągłą: \\ 
{\color{blue} plot 'dane.dat' smooth csplines}
\end{itemize}
}


\begin{frame}{Wykresy -- opis }

\begin{minipage}[t]{5cm}
Często (w zasadzie zawsze) konieczne jest umieszczenie oznaczeń osi i tytułu wykresu. 
xoff, yoff - przesunięcie o wielokrotność szerokości liter\\
Podobnie instrukcje dla osi y
\end{minipage}
\begin{minipage}[t]{5cm}
\begin{figure}
	\centering
	\includegraphics[width=5cm]{przyklad_5.eps}
%	\caption{przyklad 5}
\end{figure}
\end{minipage}

Pomocne są tu instrukcje: \\
 {\color{blue} set xlabel \{"`<label>"'\} \{<xoff>\}\{,<yoff>\} \{"`<font>\{,<size>\}"'\} \\
      show xlabel \\
}

\end{frame}

\begin{frame}{Wykresy -- opis}

Tytuły \\
{\color{blue} set title \{"`<title-text>"'\} \{<xoff>\}\{,<yoff>\} \{"`<font>,\{<size>\}"'\} \\
}
\begin{minipage}[t]{5cm}
Musi być aktywna opcja pokazująca legendę: \\
{\color{blue} set key} \\
Instrukcje tą można skrócić do samego "`t"'. \\
Umieszczamy ją w linii generującej wykres:
\end{minipage}
\begin{minipage}[t]{5cm}
\begin{figure}
	\centering
	\includegraphics[width=5cm]{przyklad_6.eps}
%	\caption{przyklad 6}
\end{figure}
\end{minipage}


{\color{blue} p sin(x) t "`wykres funkcji"' w l lt 6 lw 2} \\
\end{frame}

\begin{frame}{Wykresy -- tytuły}
Dodatkowe instrukcje dotyczące tytułów wykresów:\\
{\color{blue}
set key \{  left | right | top | bottom | outside | below \\
 \hspace{1cm}              | <position>\} \\
 \hspace{1cm}             \{Left | Right\} \{\{no\}reverse\} \\
 \hspace{1cm}             \{samplen <sample length>\} \{spacing <vertical spacing>\} \\
 \hspace{1cm}             \{width <width increment>\} \\
 \hspace{1cm}             \{title "<text>"\} \\
 \hspace{1cm}             \{\{no\}box \{ \{linestyle | ls <line style>\} \\
 \hspace{1cm}                        | \{linetype | lt <line type>\} \\
 \hspace{2cm}                          \{linewidth | lw <line width>\}\}\} \\
      set nokey \\
      show key \\
}
\end{frame}

\begin{frame}{Wykresy -- tytuł}
\onslide<1->{
Położenie tekstu na stronie: left, right, top, bottom, outside i below \\
Położenie wewnątrz pudełka: Left lub Right \\
Domyślne ustawienia: \\
right, top, Right, noreverse, samplen 4, spacing 1.25, title "", nobox \\}
\vspace{0.5cm}
\onslide<2>
Podobnie możemy zastosować do etykiet: \\
{\color{blue}  set label \{<tag>\} \{"`<label text>"'\} \{at <position>\} \\
 \hspace{1cm}               \{<justification>\} \{\{no\}rotate\} \{font "`<name><,size>"'\} \\
\hspace{1cm}      set nolabel \{<tag>\} \\
\hspace{1cm}      show label \\
}
Położenie można zdefiniować we współrzędnych wykresu {\color{blue} first}, ale znacznie wygodniej zrobić to względem wykresu {\color{blue} graph} \\
{\color{blue} set label "`y=x"' at 1,2 } lub \\ 
{\color{blue} set label "`S"' at graph 0.5,0.5 center font "`Symbol,24"'}
\end{frame}

\frame{{Znaczniki osi}
Główne znaczniki osi:\\
{\color{blue}  set xtics 5 (... ,-5, 0, 5, ...) \\
set ytics 0 2 10 (0, 2, 4, 6, 8, 10) \\
set log x; set xtics 1,100,1e6 \\
$(1, 10^2, 10^4, 10^6)$
}

Znaczniki pomocnicze (10 przedziałów, 9 znaczników): \\
{\color{blue}  set mxtics 10 }

Zastępowanie znaczników: \\
{\color{blue} set xtics ("`pon"' 1, "`wt"' 2, "`sr"' 3, "`czw"' 4) }
}
\frame{{Znaczniki osi}
Zmiana formatu znaczników:
{\color{blue} set format y "`$10\hat{}${ \%L}"'} w zapisie potęgowym \\

Ogólnie {\color{blue} set format} + nazwa osi (X, Y, Z, XY, X2, Y2) + definicja formatu\\
\% całkowita długość.dokładność \\ 
typ ('f', 'e', 'E', 'g', 'x', 'X', 'o', 't', 'l', 's', 'T', 'L', 'S', 'c', 'P') \\
domyślnie przyjęty jest 'g'.\\
't', 'l', 'T', 'L' -- stosowane są w skalach logarytmicznych, 
}
\frame{{Dodatkowe strzałki i linie}
Do umieszczania dodatkowych strzałek i linii służy polecenie {\color{blue} set arrow}.
\begin{itemize}
 \item Strzałka nr 1 od punktu (1, 2) do (1, 3): \\
{\color{blue} set arrow 1 from 1,2 to 1,3 }
\item Linia nr 3, gruba na 2pt, od (1, 2.5) do (1, 3): \\
{\color{blue} set arrow 3 from 1,2.5 to 1,3 nohead lw 2 }
 \item Usunięcie wszystkich strzałek i linii: \\
{\color{blue} unset arrow }
\item Usunięcie strzałki nr 3:
{\color{blue} unset arrow 3 }
\item Wyświetlenie wszystkich strzałek i linii: \\
{\color{blue} show arrow }
\end{itemize}

}

\begin{frame}{ Wykresy -- kilka na jednej stronie}
Kilka wykresów  \\
{\color{blue} set multiplot \\
      set nomultiplot} \\
Przykład: \\
{ \color{blue} 
      set multiplot \\
      set size 0.4,0.4 \\
      set origin 0.1,0.1 \\
      plot sin(x) \\
      set origin 0.5,0.5 \\
      plot cos(x) \\
      set nomultiplot \\
}
\end{frame}
\begin{frame}{Wykresy -- kilka}
\begin{figure}
	\centering
	\includegraphics[width=10cm]{przyklad_7.eps}
%	\caption{przyklad 6}
\end{figure}
\end{frame}
\begin{frame}{Wykresy -- automatyzacja}

{\color{blue} 
\#! /usr/bin/gnuplot \\
reset \\
set term postscript eps color \\
set output "`przyklad 7.eps"' \\
\vspace{0.5cm}
\noindent       set multiplot \\
      set size 0.4,0.4 \\
      set origin 0.1,0.1 \\
      plot sin(x) \\
      set origin 0.5,0.5  \\
      plot cos(x) \\
      set nomultiplot  \\ }
\end{frame}

\frame{
\frametitle{Wykresy -- plik wyjściowy }

Generowanie wykresów w postaci plików wymaga określenia następujących parametrów: \\
\begin{itemize}
\item nazwy pliku
\item rodzaju pliku.
\end{itemize}
\begin{block}{Uwaga}
{\large Gnuplot nie rozpoznaje typu pliku po jego rozszerzeniu! } \\
Sami musimy zadbać o zgodność rozszerzenia z zawartością!
\end{block}
}

\subsection{Zapis wyników}
\frame{
\frametitle{Wykresy -- plik wejściowy}
Wszystkie instrukcje ustawiające parametry pracy gnuplota są wprowadzane komendą: \\
 \alert{ set ''nazwa zmiennej i parametry''} \\
Nazwę pliku wyjściowego definiujemy następująco: \\
\alert{ set output ''nazwa.roz''} \\
Rodzaj i format wyjścia definiujemy następująco (najistotniejsze): \\
\alert{ epslatex, latex, linux, pdf, png, postscript, \\ pslatex, pstex, pstricks, windows, x11.}

}

\frame{
\frametitle{ Postscript}
Dostępne opcje:\\
\alert{
 set terminal postscript \{<mode>\} \{enhanced | noenhanced\} \\
  \hspace{3cm}                            \{color | monochrome\} \{solid | dashed\} \\
   \hspace{3cm}                           \{<duplexing>\} \\
  \hspace{3cm}                           \{''<fontname>''\} \{<fontsize>\} \\
}

<mode> przyjmuje wartości \alert{ landscape, portrait, eps, default} \\
<duplexing> \alert{ defaultplex, simplex or duplex} - druk dwustronny \\
''<fontname>'' nazwa czcionki; <fontsize> rozmiar \\
\alert{ default} = \alert{ landscape, monochrome, dashed, defaultplex, noenhanced, "Helvetica", 14pt.} \\
eps (Encapsulated PostScript) -  PostScript + dodatkowe informacje pozwalające na łatwiejsze importowanie plików. \\
W trybie {\bf eps} cały wykres jest zmniejszany dwukrotnie.

}

\frame{{Zapis pliku}
Zapisanie zmiennych, funkcji, bieżących
parametrów rysunku i {\bf \alert{ostatniego}} polecenia plot: \\
{\color{blue} save 'nazwa\_pliku'. }

Zapis zmiennych / funkcji / parametrów: \\
{\color{blue} save functions 'nazwa\_pliku' \\
save var 'nazwa\_pliku' \\
save set 'nazwa\_pliku' }

Odczytanie tego pliku: \\
{\color{blue} load 'nazwa\_pliku' }

Zapisanie historii sesji: \\
{\color{blue} history 'nazwa\_pliku' }
}

\section{Wykresy 3D}
\frame{
\frametitle{Wykresy 3 wymiarowe}
Wykresy 3 wymiarowe tworzy się w sposób zbliżony do 2-wym.\\
\begin{block}{Instrukcja}
 splot \{ zakres \} \\
                 funkcja | '' dane '' \{ parametry danych \}  \\
                 \{ tytuł \} \{with styl\}
                 \{, \{definitions,\} <function> ...\}
\end{block}
\begin{block}{Dopasowanie rzutu}
set view \{ <rot\_x>\{,\{<rot\_z>\}\{,\{<scale>\}\{,<scale\_z>\}\}\} | map \} \\
set view 60, 30, 1, 1 \\
           set view ,,0.5
\end{block}


}
\frame{
\frametitle{Wykresy 3 wymiarowe}
\begin{block}{Nałożenie siatki}
set contour \{base | surface | both\} \\
           unset contour \\
           base -- linie są nakładane na płaszczyznę XY, surface -- na wykresie.
\end{block}
\begin{example}
unset surface \\
           set contour \\
           set cntrparam ... \\
           set term table \\
           set out 'filename' \\
           splot ... \\
           set output \\
          set term <whatever> \\
           plot 'filename' \\
\end{example}
}
\frame{
\frametitle{Wykresy 3 wymiarowe+ 4 wymiar}
\begin{block}{Instrukcja}
set pm3d \\
           set pm3d \{ \\
                      \{ at <bst combination> \} \\
                      \{ scansautomatic | scansforward | scansbackward \} \\
                      \{ flush \{ begin | center | end \} \} \\
                      \{ ftriangles | noftriangles \} \\
                      \{ clip1in | clip4in \} \\
                      \{ corners2color \{ mean|geomean|median|c1|c2|c3|c4 \} \} \\
                      \{ hidden3d <linestyle> | nohidden3d \} \\
                      \{ implicit | explicit \} \\
                      \{ map \} \\
                    \} \\
           show pm3d \\
           unset pm3d \\
\end{block}
}
\frame{
\frametitle{Wykresy 3 wymiarowe}
\begin{example}
set pm3d at s hidden3d 100 \\
           set style line 100 lt 5 lw 0.5 \\
           unset hidden3d \\
           unset surf \\
           splot x*x+y*y \\
\end{example}
}

\frame{
\frametitle{Wykresy 3 wymiarowe}
{
\frametitle{pm3d - przykład}

\includegraphics[width=8cm]{przykl_11.eps}
}
}
\frame{{Struktura danych}
W przypadku wykresów 3D należy wskazać  podział pomiędzy liniami $\rightarrow$ pusta linia w danych.

Jeżeli nie mają być rysowane poziome linie siatki to należy umieścić dwie puste linie pomiędzy danymi.

Do wykresu macierzy można wykorzystać format "`{\color{blue} matrix}"' \\
trzy pierwsze kolumny są interpretowane jako współrzędne, a czwarta jest przedstawiana kolorem.

}

% \section{Dopasowanie funkcji}

% \frame{{Fitowanie}

% Załóżmy, że mamy zbiór danych $\{ (x_i,y_i,\varepsilon_i )\} $, $ i = 1,...,N$, gdzie $\varepsilon_i$ jest niepewnością
% wartości $y_i$, oraz funkcja $f(x;a_1,...,a_k)$. \\ 
% Celem polecenia {\bf fit} jest minimalizacja
% $$
% \chi^2(\overline{a})=\sum_{i=1}^N \left[ \frac{y_i-f(x_i;\overline{a})}{\varepsilon_i} \right]^2
% $$
% względem $\overline{a} = a_1,...,a_k$.
% }

% \frame{
% \frametitle{Dopasowanie funkcji}
% Gnuplot oferuje moduł dopasowujący funkcje do danych!\\
% Składnia:\\
% set fit \{logfile \{''<filename>''\}\} \{\{no\}errorvariables\} \\
%            unset fit \\
%            show fit \\
% \begin{block}{Dopasowanie funkcji}

% fit \{[xrange] \{[yrange]\}\} <function> '<datafile>' \\
%                \{datafile-modifiers\}
%                via '<parameter file>' | <var1>\{,<var2>,...\}
% \end{block}
% Moduł ten jest oparty na metodzie najmniejszych kwadratów nieliniowego algorytmu Marquardt-Levenberga \\
% {\it http://pl.wikipedia.org/wiki/Algorytm\_Levenberga-Marquardta}
% }
% \frame{
% \frametitle{Dopasowanie funkcji}

% \begin{example}
% set fit errorvariables \\
%             fit f(x) 'datafile' using 1:2 via a, b \\
%            print ''error of a is:'', a\_err \\
%             set label 'a=\% 6.2f', a, '+/- \% 6.2f', a\_err \\
%             plot 'datafile' using 1:2, f(x) \\
% \end{example}
% }
% \frame{
% \frametitle{Dopasowanie funkcji}
% \begin{example}
% f(x) = a*x**2 + b*x + c \\
%            g(x,y) = a*x**2 + b*y**2 + c*x*y \\
%            FIT\_LIMIT = 1e-6 \\
%            fit f(x) 'measured.dat' via 'start.par' \\
%            fit f(x) 'measured.dat' using 3:(\$7-5) via 'start.par' \\
%            fit f(x) './data/trash.dat' using 1:2:3 via a, b, c \\
%            fit g(x,y) 'surface.dat' using 1:2:3:(1) via a, b, c \\
% \end{example}
% }
% \frame{{Składnia}

% W przypadku funkcji {\color{blue} \bf fit} konieczna jest specyfikacja czy dane zawierają informację o błędzie. 

% Wynika to z faktu, że jako zmienna 'x' może być wykorzystany numer wiersza. Tym samym konieczne jest podanie czy błąd jest zawarty w danych. Przyjmuje się, że błąd dotyczy przedostaniej zmiennej 'z'.

% Ze względu na wymóg kompatybilności z poprzednimi wersjami słowo kluczowe {\color{blue} \bf 'error'} nie jest wymagane.

% Domyślna kolejność zmiennych:
% \begin{tabular}{ll}
% z  & \# 1 niezależna zmienna (nr linii) \\
% x:z & \# 1 niezależna zmienna (1 kolumna) \\
% x:z:s & \# 1 niezależna zmienna \\
% x:y:z:s & \# 2 niezależne zmienne \\
% x1:x2:x3:z:s & \# 3 niezależne zmienne \\
% x1:x2:x3:...:xN:z:s & \# N niezależnych zmiennych 
% \end{tabular}
% }

% \frame{

% \begin{example}
% \tt
% f(x) = a*x**2 + b*x + c \\
% set fit limit 1e-6 \\
% fit f(x) 'measured.dat' via a, b, c \\
% fit f(x) 'measured.dat' using 1:2:3 yerror via a, b, c \\
% fit f(x) 'measured.dat' using 1:2 noerror via a, b, c \\
% \end{example}

% }

% \frame{{Istotne informacje}
% \begin{description}
%  \item[WSSR] Weighted sum of squares of residuals = $\chi^2$
% \item[FIT\_LIMIT] dokładność dopasowania
% \item[FIT\_MAXITER] ilość iteracji
%  \end{description}

% }
% \frame{{Ocena jakości dopasowania}
% \begin{description}
%  \item[ndf] ilość punktów pomiarowych (stopni swobody),
% \item[stdfit]
% \begin{itemize}
%  \item jeśli dane miały oszacowany błąd pomiaru, nasza teoria jest poprawna i zachodzą inne warunki stosowalności polecenia fit,
% to wartość stdfit powinna być bliska 1
% \item Jeśli nie znamy niepewności pomiaru, to stdfit daje wyobrażenie o średniej wartości tego błędu (o ile teoria jest poprawna etc.)
% \end{itemize}
% \item[Macierz korelacji] w przypadku idealnym wyrazy pozadiagonalne znikają
% \item[Wyniki] jest podawany z oszacowanym błędem bezwzględnym i względnym.

%  \end{description}
% }
% \frame{{Możliwe trudności}
% \begin{itemize}
%  \item Zła teoria (nie ta funkcja);
% \item Złe wartości początkowe parametrów;
% \item Nieprawidłowe wartości niepewności pomiarowych;
% \item Silne korelacje parametrów np. $f(x)=a\exp(\frac{x-b}{c})$;
% \item Poszukiwane parametry mają różne rzędy wielkości.
% \end{itemize}
% }
% \frame{{Przykład}
% \begin{center}
%  \includegraphics[scale=0.5]{fit.eps}
% \end{center}
% % fit.eps: 0x0 pixel, 300dpi, 0.00x0.00 cm, bb=50 50 410 302
% }
% \section{Pętle}
% \frame{{Pętle}
% Od wersji 4.6 istnieje możliwość zastosowania pętli i prostych iteracji oraz intrukcji warunkowych wewnątrz poleceń {\color{blue} \bf plot} i {\color{blue} \bf set}.

% Instrukcje te występują jako: {\color{blue} \bf if, while, do, for}.

% \vspace{1cm}

% Instrukcja {\color{blue} \bf for} występuje w dwóch wersjach:

% \texttt{
% plot for [<variable> = <start> : <end> {:<increment>}] \\
% plot for [<variable> in ''string of words'']
% }
% }

% \frame{{Przykłady}

% Typowa iteracja\\
% Rysujemy funkcję sin(x) dla różnych częstości:\\
% \texttt{ plot for [i=1:3] sin(i*x)}

% \vspace{0.5cm}

% Na wzór Basha pętla bprzebiega po liście parametrów: \\
% \texttt{plot for [{\bf pliki} in "praca energia"] {\bf pliki}."dat" title {\bf pliki}}

% \vspace{0.5cm}

% {\bf Uwaga}

% Iteracja kończy się na przecinku albo końcu komendy.

% Np. \texttt{plot for [i=1:3] j=i, sin(j*x)}\\
% Narysuje jedną funkcję dla i=3.
% }
% %\end{document}
% \section{Wtyczka c++}
% \lstset{language=C, showstringspaces=false,
%   basicstyle=\footnotesize,
%   keywordstyle=\bfseries\color{green!40!black},
%   commentstyle=\itshape\color{purple!40!black},
%   identifierstyle=\color{blue},
%   stringstyle=\color{orange}}
% \begin{frame}[fragile]
% \begin{lstlisting}
% #include <iostream>
% #include <fstream>
% #include <sstream>
% #include <vector>
% #include "gnuplotpp.h"

% using namespace std;
% using namespace gnuplotpp;

% int main( void ){

%  singleplot g;

%  g.set("term png" );
%  g.set() << gpset( "output", "\"demo1.png\"" ) << endl;

%  g.x.label( "x axis" ).format( "%03.2f" ).tics(0.2).mtics(.05);
%  g.y.label( "y axis" );

% \end{lstlisting}
% \end{frame}
% \begin{frame}[fragile]
% \begin{lstlisting}


%  g.label( "single plot example", 
%                          coord( "graph 0.5", "graph 0.5" ) );

%  g.arrow( coord( "graph 0.5", "graph 0.5" ), 
%                                coord( "graph 1", "graph 1" ) );

%  g.plotfile( "data.dat", 1, 2 )<< " t 'data'" << " pt 1" ;
  
%  g.plotfunc( "x" ) ;

%  g.save( "demo1.gp" ).exec();
% \end{lstlisting}
% \end{frame}

% \frame{ 
% \begin{figure}
%  \centering
%  \includegraphics[scale=0.5]{./demo1.eps}
%  % demo3.png: 640x480 pixel, 72dpi, 22.58x16.93 cm, bb=0 0 640 480
% \end{figure}
% }

% \begin{frame}[fragile]
% \begin{lstlisting}
% #include <iostream>
% #include <fstream>
% #include <sstream>
% #include <vector>
% #include "gnuplotpp.h"

% using namespace std;
% using namespace gnuplotpp;

% int main( void ){

% 	multiplot g(3);

% 	g.setsize(1.5,1);

% 	// one example for set command -- stream style
% 	//   line need to be terminated
% 	g.buf() << gpset( "term", "png" ) << endl; 

% 	// another example for set command -- function style
% 	//   no temination of the line needed
% 	g.set( "output \"demo2.png\"" );

% \end{lstlisting}
% \end{frame}
% \begin{frame}[fragile]
% \begin{lstlisting}
% 	// axis formatting can be done as this
% 	g.x.range( -5, 5 ).format( "%g" ).tics(1).label("x label");

% 	g.set( "multiplot" );

% 	// this is global label that is persistent to all graphs
% 	g.label( "This is global label",  coord( "graph 0.5", "0.1" ) );

% 	// this is local label 
% 	multiplot::iterator it=g.begin();
% 	for( ; it!=g.end() ; it++ )
% 		it->label() <<  "\"plot " << it-g.begin() << "\" at " 
%                 << coord( "graph 0.5", "0.5" ) << endl;

% 	g[0].x.format( "%4g" );
% 	g[0].x.range( 0., 1. );
% 	//	cerr << "g[0].str(): "<< g[0].buf().str() << endl;
% \end{lstlisting}
% \end{frame}

% \begin{frame}[fragile]
% \begin{lstlisting}
 
% 	for( it=g.begin(); it!=g.end() ; it++ ){
% 		it->buf() << gpset( "origin", coord(0.5*(it-g.begin()), 0 ) ) 
% 							  << endl;
% 		it->buf() << gpset( "size", coord(0.5,1) ) << endl;
% 		it->plotfunc() <<  "sin(" << it-g.begin() << "*x)+" 
% 									 << "cos(" << it-g.begin()+2 << "*x)";
% 	}

% 	g.postset() << gpset( "nomulti" ) << endl;

% 	g.save( "demo2.gp" ).exec();
 
% }

% \end{lstlisting}
% \end{frame}

% \frame{ 
% \begin{figure}
%  \centering
%  \includegraphics[scale=0.3]{./demo2.eps}
%  % demo3.png: 640x480 pixel, 72dpi, 22.58x16.93 cm, bb=0 0 640 480
% \end{figure}
% }


% \begin{frame}[fragile]
% \begin{lstlisting}
% #include <iostream>
% #include <fstream>
% #include <sstream>
% #include <vector>
% #include "gnuplotpp.h"

% using namespace std;
% using namespace gnuplotpp;

% int main( void ){

%  gptable g(3,3);

%  // cerr << g.set().str() << endl;
%  // exit(1);

%  g.set() << gpset( "term", "png" ) << endl 
% 	 << gpset( "output", "\"demo3.png\"" ) << endl;

% \end{lstlisting}
% \end{frame}
% \begin{frame}[fragile]
% \begin{lstlisting}

%  g.setsize( 1,1 );

%  g.x.range( -5, 5 );

%  for( int i=0; i<g.xsize(); i++ )
%    for( int j=0; j<g.ysize(); j++ ){
%       g[i][j].plotfunc() <<  "sin(" << i+1 << "*x)+" 
% 				<< "cos(" << j+1 << "*x)";
%    }

%  // cerr << g << "######" << endl;

%   g.save( "demo3.gp" ).exec();
 
% }
% \end{lstlisting}
% \end{frame}

% \frame{ 
% \begin{figure}
%  \centering
%  \includegraphics[scale=0.5]{./demo3.eps}
%  % demo3.png: 640x480 pixel, 72dpi, 22.58x16.93 cm, bb=0 0 640 480
% \end{figure}
% }

\end{document}


